\textbf{1. Kiến thức cần nhớ}
\begin{itemize}
\item Gọi $d$ là khoảng cách từ tâm $O$ đến đường thẳng $\Delta$, $R$ là bán kính đường tròn $(O)$:
\begin{itemize}
\item $\Delta$ cắt $(O) \iff \boxed{d < R}$ (có $2$ điểm chung).
\item $\Delta$ tiếp xúc $(O) \iff \boxed{d = R}$ (có $1$ điểm chung duy nhất, $\Delta$ gọi là \textbf{tiếp tuyến}, điểm chung gọi là \textbf{tiếp điểm}).
\item $\Delta$ không giao $(O) \iff \boxed{d > R}$ (không có điểm chung).
\end{itemize}
\item \textbf{Tính chất của tiếp tuyến}: Tiếp tuyến của đường tròn \textbf{vuông góc với bán kính} đi qua tiếp điểm:
$$\boxed{\Delta\text{ là tiếp tuyến tại } A \implies \Delta \perp OA\text{ tại } A}$$
\item \textbf{Dấu hiệu nhận biết tiếp tuyến}: Nếu một đường thẳng đi qua một điểm của đường tròn và vuông góc với bán kính đi qua điểm đó thì đường thẳng đó là tiếp tuyến của đường tròn.
\item \textbf{Tính chất hai tiếp tuyến cắt nhau}: Nếu hai tiếp tuyến của $(O)$ cắt nhau tại $M$ (tiếp điểm $A$ và $B$):
\begin{itemize}
\item $MA = MB$.
\item $MO$ là tia phân giác của góc $\widehat{AMB}$.
\item $OM$ là tia phân giác của góc $\widehat{AOB}$.
\item $MO$ là đường trung trực của đoạn thẳng $AB$ ($MO \perp AB$ tại trung điểm của $AB$).
\end{itemize}
\end{itemize}

\textbf{2. Phương pháp giải các dạng}

%% DANG 9H5?16-1
\textbf{Dạng 1. Xác định vị trí tương đối của đường thẳng và đường tròn}
\begin{enumerate}
\item Kẻ $OH \perp \Delta$, tính độ dài đoạn vuông góc $OH = d$.
\item So sánh $d$ với bán kính $R$ để kết luận số giao điểm.
\end{enumerate}

%% DANG 9H5?16-2
\textbf{Dạng 2. Tiếp tuyến: dấu hiệu nhận biết, chứng minh tiếp tuyến}
\begin{enumerate}
\item \textbf{Cách 1 (Điểm đã thuộc đường tròn)}: Chứng minh đường thẳng vuông góc với bán kính tại tiếp điểm: $\Delta \perp OA$ tại $A \in (O)$.
\item \textbf{Cách 2 (Điểm chưa thuộc đường tròn)}: Kẻ khoảng cách từ tâm $O$ đến đường thẳng, chứng minh khoảng cách đó bằng bán kính $R$ ($d = R$).
\end{enumerate}

%% DANG 9H5?16-3
\textbf{Dạng 3. Tính chất hai tiếp tuyến cắt nhau}
\begin{enumerate}
\item Vận dụng $MA = MB$ để chứng minh tam giác $MAB$ cân tại $M$.
\item Vận dụng $MO \perp AB$ và tam giác $OAM$ vuông tại $A$ để dùng hệ thức lượng:
$$OA^2 = OH \cdot OM \implies \boxed{R^2 = OH \cdot OM}$$
\item Vận dụng tính chất tia phân giác để chứng minh các góc bằng nhau hoặc tính số đo góc.
\end{enumerate}
